\subsection{Exceptional decomposition groups and Frobenius classes}

\begin{claim}{AR.8}{Known, checked; the exact Sa\"idi--Tamagawa statement}
Let $C/k$ be proper hyperbolic over a finite field, and put
$\Sigma=\{\ell:\ell\ne p\}$.  On the universal geometrically
pro-$\Sigma$ covering, let $\widetilde E_C$ consist of points whose
decomposition subgroup is commensurate with that of another point;
let $E_C$ be its image in the set of closed points of $C$.
The following are the statements supplied by the two papers.
\begin{enumerate}
\item In \cite[Definition~1.5, Lemma~1.7, Proposition~1.8(iii)--(vii),
Definition~1.10 and Corollary~1.11]{AR_ST2009}, the commensurability
map is injective above a fixed geometric point, is injective above a
rational point, has finite fibres, and has finite exceptional image
$E_C$ when $\Sigma$ is cofinite.  Decomposition subgroups of closed
points are recovered injectively outside $E_C$.
\item For a $\Sigma$-covering $f:C'\to C$,
\cite[Proposition~2.11]{AR_ST2017} proves
$E_{C'}=f^{-1}(E_C)$ if $f$ is \'etale.  If $f$ is tame Galois,
it proves $E_{C'}\subset f^{-1}(E_C)\cup R_f$, where $R_f$ is the
non-\'etale locus.
\item With $J_C$-largeness as defined in that paper,
\cite[Proposition~2.12]{AR_ST2017} places the exceptional points over
a proper subfield $F\subset\bar k$, finite of degree at most two over
$F_0=\bar k^{\ker\rho_{J_C,\Sigma'}}$.  The degree is at most one
when $C$ is nonhyperelliptic.  Proposition~2.13 gives a prime $\ell$
and finiteness of the exceptional points rational over each finite
extension of the unique $\mathbf Z_\ell$-extension of $k$; the
precise hypothesis is almost $\Sigma$-separatedness, with the stated
$J_C$-large refinements on the choice of $\ell$.
\end{enumerate}
These results do not assert $E_C=\varnothing$.
\end{claim}


\begin{proof}
We use their commensurability equivalence relation, rather than replace
it by equality of arbitrary chosen subgroup representatives.  The
stated hypotheses apply because the complement of our $\Sigma$ is
the single prime $p$.  For the first item the crucial centralizer
statement, their Lemma~1.7, applies to a closed subgroup $Z$ whose
projection to $G_k$ is injective and open: projection embeds the entire
commensurator of $Z$ in $G_k$.  Its intersection with the geometric
kernel is consequently trivial.  This proves injectivity above a
fixed geometric point.  The finite-fibre statement above a closed
point follows by its finitely many geometric conjugates.  The proof
of global finiteness first passes to a nonhyperelliptic cover, then
uses finiteness of the appropriate torsion intersection with the
difference variety in its Jacobian.  It bounds the exceptional set;
it does not show that this intersection is empty.

For the second item the universal covering is the same after passing
to an open subgroup.  Commensurability is unchanged by intersecting
with it, giving equality of the exceptional inverse images in the
\'etale case.  The tame argument allows the stated ramification locus,
which is why equality is not claimed there.

For the third, the proof of Proposition~2.12 uses
$\delta:C\times C\setminus\Delta\to J_C$,
$(x,y)\mapsto[x-y]$.  Its geometric fibres have cardinality at most
$e$, with $e=2$ in the hyperelliptic case and $e=1$ otherwise.
An exceptional pair has difference in the $\Sigma'$-primary torsion,
which is rational over $F_0$.  Thus both coordinates are rational
over a subextension of degree at most $e$.  This is the actual
descent comparison producing $F$.  Almost separation means exactly
that the exceptional residue fields are contained in such a proper
$F$.  Choosing a nonzero $\mathbf Z_\ell$ projection of its Galois
subgroup gives Proposition~2.13's finite-intersection property and
the asserted finiteness.  None of these comparisons removes the
possible exceptional points.
\end{proof}

\begin{claim}{AR.9}{Proved; A1 in the exceptional-set language}
For a fixed proper hyperbolic $C/k$, absence of collisions after
\emph{every} finite constant extension is equivalent to $E_C$ being
empty.  Universal A1 is therefore equivalent to emptiness of all these
exceptional sets.  Absence of collisions only among $C(k)$ is the
weaker assertion that no two distinct rational-point decomposition
sections become geometrically conjugate at that field.
By AR.3 all these exceptional sets are empty for
$\Sigma=\{\ell:\ell\ne p\}$.  This last conclusion is a consequence
of the proof in the present return, not an assertion attributed to
either of the cited papers.
\end{claim}
\begin{proof}
After choosing compatible lifts, conjugate rational-point sections give
equal decomposition subgroups at two different geometric points.
They are commensurate, hence exceptional.  Conversely, take an
exceptional pair upstairs.  Their geometric images are distinct by
AR.8.  Their decomposition subgroups have open intersection, because
they are commensurate.  Enlarge the finite constant field so that both
geometric points are rational and the arithmetic quotient of this
intersection contains the new $G_L$.  Each decomposition group
projects injectively for a proper curve.  Restricting the two inverse
sections to $G_L$ therefore gives the same subgroup and the same
section.  These are sections of distinct $L$-points.  This comparison
also explains why the single-field and all-extension statements must
be kept separate.

The ``Frobenius'' terminology of L7 is exact with this qualification:
$G_k=\widehat{\mathbf Z}$ has a distinguished arithmetic generator,
and a continuous section is determined by its value on that generator.
Geometric conjugacy of these elements is geometric conjugacy of the
sections.  Recovering the curve from an outer isomorphism of its
fundamental group supplies no extra argument for excluding the
exceptional pair just described.  The full abelian class-field map
does separate points, but its prime-to-$p$ quotient leaves precisely
the possible $p$-primary difference in AR.1.
Applying AR.3 after the finite constant extension used above excludes
every exceptional pair and proves the asserted emptiness.
\end{proof}

\subsection{The field-faithfulness margin of scan 36}

\begin{claim}{AR.10}{Proved; the geometrically prime-to-$p$ field functor is faithful}
Let $F,K$ be fields finitely generated over $\mathbf F_p$.  Let
$E_F^{(p')}$ and $E_K^{(p')}$ denote their absolute Galois groupoids
with the kernel over $\overline{\mathbf F}_p$ replaced by its maximal
prime-to-$p$ quotient, retaining the maps to $BG_{\mathbf F_p}$.
If two $\mathbf F_p$-embeddings $u,v:F\hookrightarrow K$ induce
isomorphic morphisms of these groupoids \emph{over}
$BG_{\mathbf F_p}$, including the specified relative commutation
isomorphisms, then $u=v$.
Thus the margin in cote~149, scan~36, concerning Th\'eor\`eme~2
(scan~22), has a correct proof using the genus-zero target
$\mathbf G_m$ and its products.  This assertion does not settle A1
for proper curves.
\end{claim}
\begin{proof}
For $n$ prime to $p$, the torsor $X^n=a$ of $a\in K^\times$ has
geometric monodromy contained in $\mu_n$, hence its Kummer class
factors through $E_K^{(p')}$.  The coefficient group $\mu_n$ is
pulled back from $BG_{\mathbf F_p}$.  We claim that
\[
 K^\times\longrightarrow
       \prod_{(n,p)=1}H^1(E_K^{(p')},\mu_n)
                                                        \tag{AR.10.1}
\]
is injective.  If $a$ is killed by every Kummer class, it is an
$n$th power in $K$ for every such $n$.  On a normal proper model
over the finite algebraic constant field $k_c$, every divisorial
valuation of $a$ is divisible by all positive integers prime to
$p$, hence is zero.  Therefore $a\in k_c^\times$.  An $n$th root
in $K$ of a nonzero constant is algebraic over $k_c$, and so lies
in $k_c$.  Taking $n=\#k_c^\times$, which is prime to $p$, forces
$a=1$.  If $K$ itself is finite, the last argument alone suffices.

For every $a\in F^\times$, functorial pullback of its Kummer
torsors under the two relative groupoid morphisms gives equal
classes in (AR.10.1).  The relative commutation isomorphism is what
identifies their coefficient systems without an uncontrolled
cyclotomic automorphism.  Thus $u(a)/v(a)=1$ by injectivity.
Both embeddings also preserve zero, proving equality.

Geometrically this test is evaluation of the coordinate of
$\mathbf G_m=\mathbf P^1\setminus\{0,\infty\}$.  To test a finite
list of nonzero generators one takes its product.  Each multiplicative
coordinate separates by the same Kummer argument; no proper-curve
section assertion has entered the proof.

Absolute Frobenius requires care.  It is isomorphic to the identity
on the absolute \'etale topos, but the isomorphism induces arithmetic
Frobenius on the base fibre functor.  It is generally not an
isomorphism between the two \emph{specified relative} morphisms.
Erasing that commutation datum, or taking the maximal prime-to-$p$
quotient of the entire arithmetic group, would change the assertion
we have proved.  Scan~22's relative-groupoid convention supports
precisely the convention retained here.
\end{proof}

\subsection{The norm module and its necessary point-origin condition}

\begin{claim}{AR.11}{Proved; construction and additivity of the norm module}
Fix $C/k$ and $P\in C(k)$.  Over $\bar k$, take all finite connected
Galois prime-to-$p$ covers $D_N\to C_{\bar k}$ indexed by open normal
subgroups $N\triangleleft\Pi_C$, with compatible lifts $P_N$.
There is a well-defined module
\[
 M_P=\varprojlim_N J_{D_N}(\bar k)[p^\infty]
                \quad\text{under norm},
 \qquad \mathbf Z_p[[\Pi_C]]\times M_P\longrightarrow M_P.
\]
Arithmetic Frobenius acts semilinearly: if $\alpha$ is its action on
$\Pi_C$, then $F(am)=\alpha(a)F(m)$.  It need not commute with the
deck action.  The module itself need not be a torsion group of bounded
exponent, although each coordinate is $p$-primary torsion.

If $P,Q,R$ have equal point-section classes, their compatible lifted
differences satisfy $T_{PR}=T_{PQ}+T_{QR}$ in this same based module.
The formula is canonical after choosing the compatible based tower;
independently based towers require their path identification first.
\end{claim}
\begin{proof}
The finite maps $D_{N'}\to D_N$, for $N'\subset N$, induce norm
homomorphisms on Jacobians.  They compose transitively and preserve
$p$-primary torsion.  Deck actions commute with the norms through the
quotient maps $\Pi_C/N'\to\Pi_C/N$.  At a fixed coordinate $m_N$
of finite order $p^a$, a completed group-ring element acts through
$(\mathbf Z/p^a)[\Pi_C/N]$.  This defines the stated action and its
compatibility, independently of the representative in the completed
ring.

The section $s_P$ acts on the entire pointed tower; it sends the
index $N$ to $\alpha(N)$ and the chosen point to the chosen point.
Galois functoriality of norm therefore defines $F$ on the inverse
limit and gives the semilinearity formula.  Each $N$ has a finite
orbit under $F$.  A power $F^m$ stabilizing $N$ makes this pointed
cover a cover over $k_m$; at that coordinate, $F^m$ is its usual
arithmetic Frobenius on the Jacobian.

AR.1 and uniqueness give compatible lifts $Q_N,R_N$ whenever the
sections collide.  In the divisor class group of $D_N$ the identity
\[
 [R_N-P_N]=[Q_N-P_N]+[R_N-Q_N]
\]
is exact.  Norm sends each term to the same term at the lower level,
so the identities give the claimed one in $M_P$.  Equality of
point-section classes supplies a fixed path.  Its geometric
centralizer ambiguity is trivial by the Lemma~1.7 comparison in
AR.8.  In particular no unrelated choice of origins is silently
added to the formula.
\end{proof}

\begin{claim}{AR.12}{Proved; what a collision gives in the module}
A collision of $P\ne Q$ gives a nonzero $F$-fixed vector
$T_{PQ}\in M_P$ whose coordinate at every level is a nonzero
point difference.  Every odd-order geometric deck group acts freely
on that coordinate.  In characteristic two this covers every level.
In odd characteristic the same conclusion for all prime-to-$p$
levels follows using the independently established hyperelliptic
case of R4.
The converse from an arbitrary free-orbit vector in $M_P$ is not
proved: the coordinate must come from actual points, compatibly.
\end{claim}
\begin{proof}
The torsion and compatibility assertions follow from AR.1.  The Abel
map of a positive-genus curve is injective on geometric points, so
none of these differences is zero.  Galois uniqueness makes the
family $F$-fixed, even at levels that Frobenius permutes.  For odd
deck groups, AC.1 proves freeness.

For completeness, suppose $p$ is odd and a nonidentity deck
transformation fixes the point difference at another level.  Its
two effective degree-two divisors in AC.1 are distinct and disjoint,
so $D_N$ is hyperelliptic.  The lifted points still have equal
sections in its full geometrically prime-to-$p$ extension, after a
finite field extension defining all the data.  The odd-characteristic
hyperelliptic case of R4, with its prime-to-$p$ correspondence checked
in AG.8 and its source audit, excludes such a collision.  Hence the
stabilizer is trivial at these levels too.  This last step uses
that established geometric theorem, rather than incorrectly apply
the odd-order argument to an involution.

None of the module construction forces a general torsion vector to
be represented by a divisor $Q_N-P_N$.  R3's small-support function
argument depends on that very property.  The abstract finite-module
example in AG.6 makes clear why a faithful or free action alone
cannot replace it in that proof.
\end{proof}

\begin{claim}{AR.13}{Proved; the rational tests as module projections}
Let $D\to C_{k_m}$ be a geometrically Galois prime-to-$p$ cover with
a chosen rational lift of $P$.  The projection of $M_P^F$ to its
Jacobian is contained in $J_D(k_m)[p^\infty]$.  Every homomorphism
$h:J_D\to B$ defined over $k_m$ with $B(k_m)[p]=0$ kills that
projection.  If $D$ is a nontrivial odd-degree L1 witness, its deck
group fixes the projection of every element of $M_P^F$.
Consequently $M_P^F$ has no element with free deck orbit at
\emph{every} level whenever there is an L1 witness for $(C,P)$
over any finite extension.

In particular this nonexistence assertion holds for every original
$\mathbf F_4$-rational point on every genus-two curve over $\mathbf F_4$,
by AR.4--AR.6. The first requested classes $203$ and $302$ already
follow from the degree-seven witnesses of AR.5 and AC.17.
\end{claim}
\begin{proof}
Choose an identification of the given pointed geometric cover with
the level $N$.  Its descent is the one induced by the power $F^m$
of the chosen point section, so an $F$-fixed family has an
$F^m$-fixed coordinate.  Fixed geometric torsion points are exactly
rational torsion points over $k_m$.  The assertion for $h$ follows
because its image is a rational $p$-primary point.  For L1,
$e=|G|^{-1}\pi^*\pi_*$ is a $k_m$-defined projector on this
$p$-primary group.  The complementary part vanishes by the relative
class-number condition, exactly as in AC.2.  Thus every projection
is $G$-invariant, including a possible zero projection.  Such a
coordinate cannot have a free orbit for nontrivial $G$.

This proof requires all the specified levels in the free-orbit
condition.  Although characteristic subgroups form a cofinal system
for constructing the inverse limit, freeness of a vector does not
in general descend under a norm to a quotient level.  We therefore
do not replace ``every level'' by ``a chosen cofinal system'' in
this assertion.  The projections used by L5 are likewise the typed
maps $h\circ\operatorname{pr}_D$ just proved to vanish; a
correspondence separating actual points adds the further geometric
information stated in AG.1.
\end{proof}

\begin{claim}{AR.14}{Open; remaining scope of the ghost-module proposal}
For every genus-two curve $C$ in characteristic two and every
rational point $P$ over a finite field of definition, $M_P^F$ has no
nonzero vector with free deck orbit at every normal level.
The implication from this assertion to A1 is proved.  A converse,
or a proof that arbitrary such vectors are compatible differences
of points, is not established.  Every original $\mathbf F_4$ pointed pair, including the first
requested classes $203$ and $302$, is settled by AR.13.
\end{claim}

\subsection{A rational exponent obstruction on a nonabelian cover}

\begin{claim}{AR.17}{Proved; small annihilators of the actual rational torsion group}
Let $L$ be a finite field of characteristic $p$, let $C/L$ be
smooth, proper and geometrically connected with $g(C)\ge2$, and
let $P\in C(L)$. Let $\pi:D\to C$ be a connected finite etale
Galois cover with constant nonabelian group $G$ of order $m$
prime to $p$, and choose $P'\in D(L)$ above $P$.
The group
\[
                 H=J_D(L)[p^\infty]
\]
is a finite $\mathbf Z[G]$-module, with the action induced by
the $L$-automorphisms of $D$. Suppose a nonzero element
$a=\sum_g a_g g\in\mathbf Z[G]$ annihilates $H$, and put
$d=\sum_g|a_g|$. If
\[
                    (d,m)=1,\qquad m>(d-1)^2,
\]
then $P$ collides with no other point of $C(L)$ in the full
geometrically prime-to-$p$ arithmetic fundamental group.

In particular, if $p^e$ is the exponent of $H$, with $p^e=1$
when $H=0$, it suffices that
\[
                         m>(p^e-1)^2.                 \tag{AR.17.1}
\]
For characteristic two, a constant pointed nonabelian cover
of degree $21$ meets this condition whenever its rational
two-primary Jacobian group has exponent at most four.
\end{claim}
\begin{proof}
A hypothetical collision gives, by AR.1 over the same field
$L$, a point $Q'\in D(L)$ above the distinct point $Q$ and
a class $T=[Q'-P']\in H$. The integral group-ring element
defines an $L$-endomorphism
$a:J_D\to J_D$, whose restriction to the finite group $H$
is zero by hypothesis. Thus $a[Q'-P']=0$. After scalar
extension to $\overline L$, all hypotheses of AG.5 are
satisfied. That theorem proves that this particular divisor
class cannot be annihilated by $a$, a contradiction.

For the stated numerical consequence take $a=p^e\cdot1$.
Its action on $H$ is zero, its coefficient norm is $p^e$,
and $(p^e,m)=1$ because $m$ is prime to $p$. This proves
(AR.17.1). For $p=2$, $m=21$ and $2^e\le4$, we have
$21>(4-1)^2=9$.

The exact field is essential. Passing to a larger field to
make a deck group constant can enlarge $H$ and its exponent,
so no such extension is made silently. The criterion only
requires an annihilator of the finite group of rational
torsion points; it does not require the image abelian variety
$a(J_D)$ to have no rational $p$-torsion. For the scalar
choice $a=p^e$, that image is the whole Jacobian. The
contradiction also uses that $T$ is a difference of points
in two distinct fibres. It does not apply to an arbitrary
vector in the norm module of AR.11.
\end{proof}

\subsection{Hereditary relative tests and the full normal-cover inverse limit}

\begin{claim}{AR.18}{Proved; a compactness consequence with exact directions}
Let $C/k$ be a smooth proper geometrically connected curve of genus
at least two over a finite field of characteristic two, and let
$P\in C(k)$. Use the full geometrically prime-to-two group and the
pointed norm module $M_P$ of AR.11. Consider the following three
conditions.
\begin{enumerate}
\item There are a finite extension $L/k$ and compatible pointed
finite etale covers $Z\to Y\to C_L$ such that both $Z$ and $Y$
are geometrically Galois of odd degree over $C$, the relative
cover $Z\to Y$ is nontrivial and Galois, and
$2\nmid h_L(Z)/h_L(Y)$.
\item The module $M_P^F$ has no nonzero element whose projection
has a free geometric deck orbit at every normal level.
\item There are a finite extension $L/k$, a pointed geometrically
Galois odd-degree cover $Y\to C_L$, and a pointed cyclic cover
$Z\to Y$ of odd prime degree with constant deck group over $L$,
such that $Z\to C_L$ has geometric Galois closure of odd degree
and $2\nmid h_L(Z)/h_L(Y)$.
\end{enumerate}
Then condition~1 implies condition~2, and condition~2 implies
condition~3. The source $Z$ in condition~3 is not asserted to be
normal over $C$.

More precisely, if condition~3 fails, then $M_P^F$ contains a
nonzero element killed by two with free deck orbit at every
normal level. Such an element is not asserted to arise from
points of $C$ or of its covers. This is a statement about the
entire inverse limit, not just finitely many quotients.
\end{claim}
\begin{proof}
For the first implication, choose the two normal levels of the
pointed geometric tower corresponding to $Z$ and $Y$. A vector
in $M_P^F$ has a coordinate in $J_Z(L)[2^\infty]$, by the
pointed descent comparison of AR.13. Norm and pullback for
$Z\to Y$ give the averaging projector for its geometric deck
group, even if that group is not constant over $L$. Its
complementary rational two-primary group is zero by the
relative class-number hypothesis; this is AC.18 and AG.19.
Consequently every element of $\operatorname{Gal}(Z_{\bar k}/Y_{\bar k})$
fixes the coordinate. This is a nontrivial subgroup of the
geometric deck group of $Z/C$, so the coordinate does not have
a free orbit. No assertion about a nonnormal level is used.

We prove the second implication by contraposition. Let
$\Pi=\pi_1(C_{\bar k},P)^{(2')}$, and let $\alpha$ be the
continuous automorphism defined by the arithmetic Frobenius
through the point section. Open normal $\alpha$-stable subgroups
$N$ are cofinal among all open normal subgroups of $\Pi$.
Indeed, a pointed finite geometric cover and its chosen point
descend over a finite constant extension, so its defining
subgroup has a finite $\alpha$-orbit; intersect that orbit.
Finite intersections also give a common refinement of any
finite list of such subgroups.

Write $D_N/k$ for the pointed model supplied by the point
section, $G_N=\Pi/N$ for its geometric deck group, and
\[
                  V_N=J_{D_N}[2](\bar k).
\]
This is a finite-dimensional $\mathbf F_2$-vector space.
Arithmetic Frobenius $F_N$ acts on it semilinearly for its
finite odd-order group $G_N$; no constant-deck assumption is
made. Each $G_N$ is solvable, by the verified odd-order theorem
used in AC.3.

Failure of condition~3 implies the hypothesis of AC.40 at this
level. To check the direction, any cyclic-prime test isolated
there has a normal intermediate base $D_N/B$ over $C$, a
source $D_N/A$ with odd geometric Galois closure, and the
compatible rational point over a finite constant extension.
If its relative class number were odd with a nonconstant
cyclic deck group, AC.20 would make that relative group constant
over its exact Frobenius orbit extension while preserving the
oddness. This would give condition~3. We do not make an
additional, uncontrolled extension to constantize the deck
group of $D_N/B$ over $C$.

For each $N$ let $S_N$ be the finite set of $v\in V_N^{F_N}$
such that the norm of $v$ to $D_N/H$ has free orbit under
$G_N/H$ for every normal $H\triangleleft G_N$. The cover
$D_N/H$ and its torsion can be considered geometrically;
$H$ need not be Frobenius-stable. Pullback on two-torsion is
injective, because the degree is odd, and
\[
 (D_N\to D_N/H)^*\operatorname{Nm}_{D_N/(D_N/H)}(v)
                         =|H|e_Hv.
\]
Its equivariance and the invertibility of $|H|$ identify the
required stabilizer with $H$. Thus AC.40 proves $S_N\ne\varnothing$.

If $N'\subset N$, norm is a map of finite sets
\[
        S_{N'}\longrightarrow S_N,
        \qquad v\longmapsto\operatorname{Nm}_{D_{N'}/D_N}(v).
\]
It commutes with arithmetic Frobenius and preserves
two-torsion. For every normal quotient of $D_N/C$, its
composite with the next norm is the norm to that same normal
quotient of $D_{N'}/C$. Hence it preserves every freeness
condition defining $S_N$.

The inverse limit of these nonempty finite sets is nonempty.
Here is the precise compactness argument, which requires no
surjectivity of the restricted norm maps. In the compact
product $\prod_N S_N$, each compatibility equation defines a
closed subset. Any finite list of equations is satisfied by
choosing one element at a common refinement and taking its
norms to all coordinates occurring in the list. The remaining
coordinates can be chosen arbitrarily. The finite-intersection
property therefore gives a compatible family $(v_N)_N$.

This family determines an element of $M_P[2]$, by cofinality,
and it is fixed by $F$ because every displayed coordinate is
fixed. Given an arbitrary normal geometric level, choose an
$\alpha$-stable refinement. Its norm to that quotient has
free deck orbit by the definition of $S_N$, so freeness holds
at every normal level, including ones that Frobenius permutes.
There is a nontrivial odd-degree geometric cover of $C$:
for any odd prime $\ell$, its geometric abelian fundamental
group has a quotient $\mathbf Z/\ell$. At a stable refinement
dominating that quotient, a free orbit cannot be the orbit
of zero. Thus the inverse-limit element is nonzero.

The first implication only treats normal sources. In the
contrapositive construction, AC.40 isolates a cover with
normal base but possibly nonnormal source. Freeness at every
normal level does not imply freeness after norm to that
nonnormal source. We therefore claim precisely the two
implications stated, and neither an equivalence of all three
conditions nor a point-origin assertion.
\end{proof}

\subsection{A high-rank abelian tower and the missing exponent estimate}

\begin{claim}{AR.19}{Proved; the elliptic translation obstruction}
In AG.5, its conclusion holds whenever the prime-to-$p$
group $G$ cannot be generated by two elements as an abelian
group. More explicitly, it holds if $G$ is nonabelian, or
if $G$ is abelian and some Sylow subgroup needs at least
three generators. Under the same group hypothesis, the
rational annihilator criterion AR.17 holds with no change
in its field, pointing, or numerical conditions.
For a scalar annihilator $p^e\cdot1$, it is enough that
the cover is geometrically Galois, with all other data
and the chosen lift defined over $L$; rational deck
transformations are unnecessary.
\end{claim}
\begin{proof}
We retain the notation of AG.5. If $a[Q'-P']=0$, the
nonzero divisor supported on the two disjoint free orbits
gives a degree-$d$ function $f$ on $D$. The field generated
by its $G$-translates defines a smooth proper curve $Z$ and
a finite map $D\to Z$ of degree $e$, with
\[
 e\mid d,\qquad g(Z)\le(d/e-1)^2.
\]
These are exactly the Castelnuovo--Severi induction and
comparison of function fields proved in AG.5. The kernel
of the $G$-action on $Z$ has order dividing $e$ and $m$,
so is trivial. A nontrivial point stabilizer would have
a prime-order subgroup partitioning a fibre divisor of
degree $e$ into equally weighted free orbits; this would
make that prime divide both $e$ and $m$. Thus the action
on $Z$ is free. The inequalities and Riemann--Hurwitz now
give $g(Z)=1$, just as in AG.5.

Every element of a free finite action on an elliptic
curve is a translation: writing it as $z\mapsto t+u(z)$,
the case $u\ne1$ would give a fixed point because $1-u$
is a surjective nonzero elliptic endomorphism. Hence $G$
embeds in the group of geometric points of an elliptic
curve. Its exponent $n$ is prime to $p$, and its image
lies in the finite group
\[
                  E[n](\overline k)\simeq
                  (\mathbf Z/n\mathbf Z)^2.
\]
A subgroup of this group is abelian and its Sylow
subgroups need at most two generators. This contradicts
the hypothesis on $G$ and proves the geometric claim.
For AR.17, a collision supplies the rational point
difference $T$, and the asserted annihilator of the
actual finite group $J_D(L)[p^\infty]$ kills $T$.
The geometric result just proved excludes that equality.
In particular the scalar annihilator $p^e$ works when
$m>(p^e-1)^2$. For this scalar case the deck group need
not be constant over $L$: the scalar kills the rational
lifted difference, and all subsequent translates and
function-field comparisons take place over the algebraic
closure. R1 supplies that lifted difference also for a
pointed cover which is only geometrically Galois. No action
of the nonrational deck transformations on $J_D(L)$ is
being asserted.
\end{proof}

\begin{claim}{AR.20}{Proved; an exact conditional abelian-tower route}
Let $C/k$ be smooth, proper and geometrically connected,
where $k=\mathbf F_q$ has characteristic $p$, let
$g(C)=g\ge2$, and let $P\in C(k)$. Fix an odd prime
$\ell\ne p$. Write $J=J_C$ and let
$j_P:C\to J$ be the Abel morphism $x\mapsto[x-P]$.
For $n\ge1$ form the fibre product over $k$
\[
 D_n=C\mathbin{\times}_{J,[\ell^n]}J,
 \qquad \pi_n:D_n\longrightarrow C.
                                                        \tag{AR.20.1}
\]
This is a geometrically connected finite etale cover
of degree $\ell^{2gn}$, pointed by $(P,0)$. Let $m_n$
be the order of arithmetic Frobenius on $J[\ell^n]$,
put $L_n=\mathbf F_{q^{m_n}}$, and let $E_n$ be the
exponent of the finite group $J_{D_n}(k)[p^\infty]$,
with exponent one for the zero group. If for one $n$
\[
                       E_n-1<\ell^{gn},                 \tag{AR.20.2}
\]
then $P$ collides with no other point of $C(k)$ in
the full geometrically prime-to-$p$ arithmetic group.
In particular an estimate $E_n=o(\ell^{gn})$ would
suffice for this $(C,P,\ell)$.

There is an integer $t\ge1$ such that
\[
           m_n=m_1\ell^{\max(0,n-t)}\quad(n\ge1).
                                                        \tag{AR.20.3}
\]
Consequently the exponents of $J(L_n)[p^\infty]$ are
uniformly bounded. This last bound concerns the base
Jacobian only; no bound for the relative parts of
$J_{D_n}$, and hence no universal instance of
(AR.20.2), is asserted here. If $E_n^{\rm const}$
is the exponent over the constant-deck field $L_n$,
then $E_n\mid E_n^{\rm const}$, so the analogous
bound over $L_n$ would also suffice, but that field
enlargement is unnecessary for the scalar criterion.
\end{claim}
\begin{proof}
Multiplication by $\ell^n$ on $J$ is finite etale,
so its pullback has the displayed degree and is etale.
The Abel morphism induces the usual isomorphism on
geometric abelian pro-$\ell$ fundamental groups.
Its monodromy onto $J[\ell^n]$ is therefore surjective,
which proves geometric connectedness. The pair $(P,0)$
is $k$-rational. Over $L_n$ the deck group is the
constant group
\[
                  G_n\simeq(\mathbf Z/\ell^n)^{2g},
\]
and the complete fibre over $P$ is rational. This
group needs $2g>2$ generators. Without extending $k$,
R1 supplies a $k$-rational lifted difference killed by
$E_n$: its hypotheses only require a geometrically
prime-to-$p$ Galois closure and the chosen rational
lift. Apply the nonconstant-deck scalar form of AR.19 with
$a=E_n\cdot1$, $d=E_n$ and $m=\ell^{2gn}$.
Their greatest common divisor is one, and (AR.20.2)
is exactly the required strict inequality.

Let $A\in\operatorname{GL}_{2g}(\mathbf Z_\ell)$
be arithmetic Frobenius on $T_\ell J$. The matrix
$A^{m_1}-1$ is nonzero: otherwise every eigenvalue
of Frobenius would be a root of unity, contrary
to its complex absolute value $q^{1/2}>1$.
Let $t$ be the minimum $\ell$-adic valuation of its
entries. Then $t\ge1$. For an odd prime $\ell$, the
binomial formula gives
\[
 v_\ell\bigl((1+\ell^t B)^{\ell^r}-1\bigr)=t+r
 \quad(r\ge0)
\]
when at least one entry of $B$ is a unit; matrix
valuation here means the minimum entry valuation.
Indeed the first term at each powering has that
valuation and every later term has strictly larger
valuation. Raising to an integer prime to $\ell$
does not change it. An exponent making $A$ the
identity modulo $\ell^n$ must first be a multiple
of $m_1$, and the same calculation then proves
(AR.20.3). In particular $v_p(m_n)=v_p(m_1)$.

For completeness we prove the asserted base bound.
On the geometric etale $p$-divisible group of $J$,
Frobenius is represented by a matrix
$B\in\operatorname{GL}_{f_p(J)}(\mathbf Z_p)$.
Its eigenvalues $\alpha_i$ are algebraic $p$-adic
units, and none is a root of unity. The rational
$p$-primary group over $\mathbf F_{q^b}$ is the
kernel of $B^b-1$ on $(\mathbf Q_p/\mathbf Z_p)^{f_p(J)}$.
Since $B^b-1$ is invertible over $\mathbf Q_p$,
its order has $p$-valuation
\[
                  \sum_i v_p(\alpha_i^b-1).             \tag{AR.20.4}
\]
Choose a finite $p$-adic splitting field. If the
residue of $\alpha_i$ has order $d_i$, then the
summand is zero unless $d_i\mid b$. Otherwise
write $b=d_i p^c u$, with $p\nmid u$. Factoring
$X^u-1$ at $X=\alpha_i^{d_i p^c}\equiv1$ shows
\[
 v_p(\alpha_i^b-1)=
 v_p(\alpha_i^{d_i p^c}-1),
\]
because the second factor reduces to the unit $u$.
For $b=m_n$ the integer $c=v_p(m_1)$ is fixed.
Thus every summand in (AR.20.4) belongs to a fixed
finite set of finite numbers. The orders, and
therefore the exponents, of the base groups are
uniformly bounded. The case $f_p(J)=0$ is immediate.

The comparison does not identify these base groups
with the entire rational torsion of $J_{D_n}$. The
new character components are the precise missing
input in the proposed route; omitting them would
invalidate the conclusion about a collision.
\end{proof}
